\section{A fifth-order congruence at upper-half primes}
\label{sec:upper-half-fifth-order}

This section concerns weight seven and the original four consecutive blocks.
It proves a local arithmetic improvement, including the rational boundary.
All congruences of rational numbers are in $\Z_{(p)}$.

\begin{theorem}\label{thm:upper-half-row}
Let $n$ be even, let $p>7$ be prime with $p<n<2p$, and let
$p\le m\le n$. Write $H_2=\sum_{h=1}^{p-1}h^{-2}$.
For the unscaled coefficients of $T_{7,n,m}$ one has
\begin{align}
 98p^6b(n,m)+153p^3a_3(n,m)&\equiv0\pmod{p^2},\label{eq:upper-half-second}\\
 174p^7b(n,m)+(191+87p^2H_2)p^4a_3(n,m)
       +247p^2a_5(n,m)&\equiv0\pmod{p^5}.\label{eq:upper-half-fifth}
\end{align}
In the first congruence both summands are $p$-integral.
For the four-block cancellation lattice, if additionally $p\le n-3$
and $p\ne29$, then $p^5\mid\Gamma_{7,n}$.
\end{theorem}

\begin{proof}
Put $n=p+a$, $m=p+b$, so $a$ is odd and $0\le b\le a<p$.
Write $c_{h,k}$ for the order-$k$ coefficient at $t=-h$.
For $0\le v\le a$ form the rational function
\[
 R_v(X)=p^7T_{7,n,m}(pX-v).
\]
Its principal parts at $X=0,-1$ have coefficients
\[
 A_{v,k}=p^{7-k}c_{v,k},\qquad
 B_{v,k}=p^{7-k}c_{p+v,k}.
\]
Let $F_v$ be the sum of just these two principal parts and put $F=\sum_vF_v$.
The following reduction is coefficientwise, including Laurent coefficients
at the two indicated poles:
\[
 F_v(X)\equiv R(X)P_v(X)\pmod{p^5},\qquad
 R(X)=\frac{(X-1)(X+2)}{X^7(X+1)^7},\quad \deg P_v\le4.       \tag{A}
\]
Here is a direct justification that also covers numerator carries.
The factorial prefactor has valuation five. The denominator has precisely
the two $p$-divisible factors $pX,p(X+1)$, each to order seven.
The numerator contains $p(X-1),p(X+2)$, and possibly $p(X-2)$ and
$p(X+3)$. Thus, after factoring out $R$, the expression is $p^c$
times a polynomial of degree $c$, times a quotient of products of
$p$-adic-unit linear factors, where $c=0,1,2$.
Expand the unit factors through order $4-c$ in $pX$.
The result has degree at most four modulo $p^5$.
Taking principal parts preserves this congruence, since the remaining
factors are units at $0,-1$ and $RP_v$ is proper with only these poles.

The poles $a<h<p$ omitted from $F$ have
\[
 p^{7-k}c_{h,k}\in p^7\Z_{(p)}.                              \tag{B}
\]
Indeed $p\mid\binom nh$, so the leading Laurent coefficient is divisible
by $p^7$. Each jet of degree $u$ has valuation at least $-u$:
all root differences have absolute value below $4p<p^2$ and each
elementary reciprocal factor loses at most one power per degree.
The jet degrees are at most six, below $p$, so their factorial
denominators introduce no further $p$-adic loss.
This proves (B) and shows that all $A_{v,k},B_{v,k}$ are $p$-integral.

Reflection of the original kernel gives $F(-1-X)=-F(X)$ exactly.
The function $R$ is reflection-even. Consequently the degree-four
polynomial $\sum_vP_v$ is reflection-odd modulo $p^5$, and hence
\[
 F\equiv A W_7+B W_6\pmod{p^5},\qquad
 W_7=(2X+1)R,\quad W_6=X(X+1)W_7,                            \tag{C}
\]
with $A,B\in\Z_{(p)}$.
Cancellation of the monic numerator and denominator in this step is
valid over $\Z/p^5\Z$, even though that ring has zero divisors.
The rational approximation $RP_v$ in (A) decays at least as $X^{-8}$.
Consequently the coefficients of $X^{-1},\ldots,X^{-7}$ in $F_v$
at infinity vanish modulo $p^5$; no such exact decay of $F_v$ is assumed.

We next verify that no harmonic boundary was discarded. Let
$H_{v,k}=\sum_{j=1}^v j^{-k}$ and $H_{p-1,k}=\sum_{j=1}^{p-1}j^{-k}$.
Splitting the multiple of $p$ from $H_{p+v,k}$ gives
\[
 H_{p+v,k}-p^{-k}
 =H_{p-1,k}+\sum_{j=1}^v(p+j)^{-k}.
\]
Expand the last summands modulo $p^{5-k}$ for $k<5$.
The coefficient multiplying $p^dH_{v,d}$ in the regular boundary error is
\[
 A_{v,d}+\sum_{k=1}^{d}(-1)^{d-k}\binom{d-1}{k-1}B_{v,k}.
\]
This is exactly the coefficient of $X^{-d}$ in $F_v$ at infinity,
and vanishes modulo $p^5$ for $1\le d\le4$ by (A).
Terms with $k\ge5$, and the omitted poles in (B), already vanish.
Writing $B_k=\sum_vB_{v,k}$ and $b_F=-\sum_k B_k$, we obtain
\[
 p^7b\equiv b_F-\sum_{k=1}^4p^kH_{p-1,k}B_k\pmod{p^5}.      \tag{D}
\]
The finite-field power-sum identity gives
$H_{p-1,2},H_{p-1,4}\in p\Z_{(p)}$.
Pairing $j$ and $p-j$ then gives
$H_{p-1,1},H_{p-1,3}\in p^2\Z_{(p)}$.
Reflection and the vanishing $X^{-1}$ coefficient give $B_1=0$ modulo
$p^5$; the vanishing $X^{-3}$ coefficient gives $B_2=B_3$ modulo $p^5$.
By (B), $2B_3\equiv p^4a_3$ modulo $p^5$. Therefore (D) becomes
\[
 p^7b\equiv b_F-\frac{p^2H_2}{2}\,p^4a_3\pmod{p^5}.        \tag{E}
\]

Direct partial fractions of the two fixed rational functions in (C) give
\[
\begin{array}{c|rrr}
 &a_3&a_5&b_F\\\hline
 W_7&-196&-64&306\\
 W_6&66&18&-98
\end{array}
\]
Both rows satisfy $174b_F+191a_3+247a_5=0$.
Apply this relation in (C), use (B), and substitute (E).
This proves \eqref{eq:upper-half-fifth}.

For \eqref{eq:upper-half-second}, repeat (A) modulo $p^3$.
Now $\deg P_v\le2$, so reflection leaves only $A W_7$.
Modulo $p$ the polynomials $P_v$ are constant; their reflection-odd sum
is zero, so $A\in p\Z_{(p)}$.
The boundary error (E) vanishes modulo $p^3$.
The relation $98\cdot306+153\cdot(-196)=0$, followed by division by $p$,
proves the first congruence and its integrality assertions.

Finally let $w$ be any integer vector cancelling the $a_3,a_5$ rows.
Apply \eqref{eq:upper-half-fifth} columnwise and pair with $w$.
Its two lower-row terms vanish exactly, giving
$p^5\mid174p^7\sum_jw_jb(n,n-j)$ in $\Z_{(p)}$.
For $p>7$, $p\ne29$, the factor $174=6\cdot29$ is a unit.
Moreover $v_p(d_n)=1$, so $p^5\mid Z_n(w)$ as an integer.
Every target-row entry is divisible by $p^7$, since $a_7$ is an integer.
Thus $p^5$ divides both coefficient images of every integer kernel vector,
proving $p^5\mid\Gamma_{7,n}$ without dividing raw minors.
\end{proof}

\begin{corollary}
Combining Theorem~\ref{thm:upper-half-row} with the earlier divisor $A_n$
gives the stronger weight-seven content estimate
\[
 \liminf_{\substack{n\to\infty\\n\text{ even}}}
 \frac{\log\Gamma_{7,n}}n
 \ge 2+\frac{\pi}{4\sqrt3}+\frac34\log3-\frac12
 =2.777409057559\ldots .
\]
\end{corollary}
\begin{proof}
Let $Q_n$ be the product of previously eligible primes
$n/2<p\le n-3$, omitting 29. They already occur to the first power
in $A_n$. The theorem supplies four additional powers, so
$A_nQ_n^4\mid\Gamma_{7,n}$.
The finitely many fixed exclusions and primes dividing the earlier
boundary polynomial have logarithmic weight $O(\log n)$.
The prime number theorem gives $\log Q_n/n\to1/2$.
Add this gain of two to the previously proved rate for $A_n$.
\end{proof}

\paragraph{Attempt to close the complete cost.}
For completeness the joint-minor saving must also survive normalization.
On the primes in $Q_n$, the two naturally scaled lower rows used in the
earlier Cramer envelope have all entries divisible by $p$.
Their two-minor gcd therefore has valuation at least two.
The new content bound gives $v_p(J_n)\ge2+5=7$, whereas the earlier
joint divisor $P_n^2A_n$ supplied three powers.
Consequently the improved joint divisor is $P_n^2A_nQ_n^4$;
this justifies the following cost substitution without counting a Smith
factor twice.
The previous elementary Cramer envelope for the normalized weight-seven
cost had upper exponent $14.995736480048\ldots$.
Substituting $A_nQ_n^4$ for $A_n$ decreases this upper exponent by two,
to $12.995736480048\ldots$.
It therefore does not prove decay of the primitive forms.
The signed polynomial moment bound remains valid, but its second integer
minimum has not been bounded at the required asymptotic scale.
The new theorem settles a higher local content estimate; it does not
establish the irrationality of $\zeta(7)$ or of every odd zeta value.

\paragraph{Verification scope.}
Exact Laurent reconstruction checks both congruences in every one of
the 1,380 nonboundary prime-index pairs through $n=240$.
The fifth-order identity has 5,520 individual column checks.
The Lean module \texttt{OddZeta/PadicRowTransfer.lean} checks the algebraic
transfer of one- and two-row congruences to kernel divisibility, together
with the fixed coarse row relation. The local rational-function reduction,
harmonic boundary argument, and prime-number-theorem corollary above are
paper proofs; their full formalization is not claimed.
